% =====================================================================
% 虚拟校园系统 · 软件设计说明书（LaTeX 版）
% 编译方式：XeLaTeX（Overleaf / Leaflet 选择 XeLaTeX 即可）
% 依赖宏包：ctex, geometry, booktabs, longtable, array, hyperref, xcolor
% =====================================================================
\documentclass[12pt,a4paper]{ctexart}
\usepackage[margin=2.5cm]{geometry}
\usepackage{booktabs}
\usepackage{longtable}
\usepackage{array}
\usepackage{graphicx}
\usepackage[colorlinks=true,linkcolor=blue]{hyperref}
\usepackage{xcolor}

\setlength{\parindent}{2em}
\linespread{1.4}

\title{\bfseries 虚拟校园系统\\[0.4em] 软件设计说明书}
\author{撰稿：李文湖（组长）等全体成员\\[0.6em] 东南大学计算机科学与工程学院}
\date{版本号：V1.0\quad 日期：2026-08-28}

\begin{document}

\maketitle
\thispagestyle{empty}

% ========================= 修改记录 =========================
\begin{center}
\begin{tabular}{l l l l}
\toprule
修改日期 & 版本号 & 修改人 & 修改内容 \\
\midrule
2026-08-28 & V1.0 & 李文湖 & 初稿：总体设计与各模块设计 \\
\midrule
 &  &  & \\
\midrule
 &  &  & \\
\bottomrule
\end{tabular}
\end{center}

\newpage
\tableofcontents
\newpage

% ========================= 1 引言 =========================
\section{引言}

\subsection{编写目的}
本说明书面向虚拟校园系统的开发者、测试者与评审教师，给出系统的总体架构、各功能模块、网络通信、多线程与数据库的详细设计，用于指导团队按 TDD 方式并行开发，并为后续编码、测试与验收提供依据。

\subsection{背景}
\begin{itemize}
  \item \textbf{软件名称}：虚拟校园系统（vCampus）。
  \item \textbf{任务提出者}：东南大学计算机科学与工程学院（专业技能实训课程）。
  \item \textbf{开发者}：6 人开发小组（组长 1 名 + 组员 5 名）。
  \item \textbf{用户}：学生、教师、管理员。
  \item \textbf{运行环境}：客户端与服务器端均运行于 Windows 桌面环境；服务器端连接 MySQL 数据库。
\end{itemize}

\subsection{定义}
\begin{longtable}{l p{10cm}}
\toprule
\textbf{术语} & \textbf{说明} \\
\midrule
\endhead
客户端（Client） & 供学生/教师/管理员使用的桌面程序，提供 Swing 图形界面并转发请求到服务器端。 \\
服务器端（Server） & 集中承载业务逻辑与数据访问的后端程序，以多线程并发服务多个客户端。 \\
公共工程（Common） & 客户端与服务器端共享的代码工程，存放实体、Message、常量与工具类。 \\
Message & 客户端与服务器端之间传输的统一消息信封，实现 \texttt{java.io.Serializable}。 \\
DAO & 数据存储与访问组件，屏蔽底层数据库差异。 \\
纵向模块 & 每个组员独立负责、跨视图/业务/数据各层端到端实现的一个功能子系统。 \\
\bottomrule
\end{longtable}

\subsection{参考资料}
\begin{itemize}
  \item 《专业技能实训课程安排（2026）》——课程要求与提交物清单。
  \item 《软件设计说明书 DEMO》——本说明书的格式模板。
  \item 《Java 编码规范》——命名、注释与格式约定。
  \item Maven、MySQL、JUnit 官方文档。
\end{itemize}

% ========================= 2 程序系统的分析 =========================
\section{程序系统的分析}

\subsection{可行性分析}
本项目为校园信息系统的教学化实现，需求明确、范围可控，团队规模与课程周期匹配，技术上基于成熟的 Java SE + Swing + Socket + MySQL 方案，具备完全可行性。

\subsection{需求分析}
系统采用客户端/服务器（C/S）架构，包含 6 个功能模块（纵向划分）：
\begin{center}
\begin{tabular}{l l l}
\toprule
\textbf{模块} & \textbf{类型} & \textbf{说明} \\
\midrule
用户管理 & 必做 & 注册/注销、登录/登出、按角色授权 \\
学生学籍管理 & 必做 & 学籍信息登记、维护、查询 \\
选课系统 & 必做 & 选课、退课、成绩查询 \\
图书馆 & 必做 & 图书检索、借书、还书 \\
商店 & 必做 & 商品浏览、购买 \\
银行 & 选做 & 余额、充值、消费流水 \\
\bottomrule
\end{tabular}
\end{center}

\subsection{开发设计环境}
\begin{center}
\begin{tabular}{l l}
\toprule
\textbf{项} & \textbf{选型} \\
\midrule
语言/编译级别 & Java，编译 level=1.7（使用 JDK 8 工具链） \\
构建工具 & Maven 3.9.x（多模块） \\
数据库 & MySQL 8.0（Connector/J 8.0.33） \\
IDE & VS Code（备选 Eclipse 2024-12） \\
GUI & Swing + Nimbus Look \& Feel \\
测试 & JUnit 5 + Mockito（TDD） \\
第三方库 & mysql-connector-j、junit-jupiter、mockito-core \\
\bottomrule
\end{tabular}
\end{center}

\subsection{程序系统的结构}
系统划分为三个 Maven 工程（模块）：
\begin{center}
\begin{tabular}{l l l}
\toprule
\textbf{工程} & \textbf{标识} & \textbf{职责} \\
\midrule
公共工程 & vcampus-common & 实体类、Message、常量、工具类（两端共享） \\
客户端 & vcampus-client & Swing 界面、客户端网络、调用远端服务 \\
服务器端 & vcampus-server & 业务服务、网络/线程池、DAO 数据访问 \\
\bottomrule
\end{tabular}
\end{center}
客户端与服务器端均依赖公共工程，通过 Socket + 对象流传输 \texttt{Message} 通信；服务器端连接 MySQL 数据库 \texttt{vCampus}。

% ========================= 3 用户管理模块 =========================
\section{用户管理模块设计说明}

\subsection{模块背景}
用于管理用户信息及用户的登录、登出、注册与授权，是其他所有业务模块的运行前提（提供统一登录态）。

\subsection{需求分析}
支持学生、教师、管理员三类角色的注册、注销、登录、登出与按角色的功能授权。

\subsection{系统设计}
采用 view / biz / vo / dao 四层结构；登录态经网络模块传递到服务器端校验后，由服务器端生成会话。

\subsection{界面设计（Client 端）}
登录窗口（用户名/密码）、注册窗口、主窗口（按角色显示可访问的功能入口）。

\subsection{模块流程图}
登录流程：\textbf{输入账号密码 → 客户端封装 Message 发送 → 服务器端校验 → 返回状态码 → 客户端跳转主窗口或提示错误}。

\subsection{类分析}
\textbf{实体类：用户（User）}
\begin{center}
\begin{tabular}{l l l l}
\toprule
序号 & 名称 & 类型 & 约束/备注 \\
\midrule
1 & id & String & 登录名（唯一） \\
2 & pwd & String & 6--16 个字符，密码 \\
3 & age & Integer & 非 0 \\
4 & sex & String & 男 / 女 \\
5 & role & String & 学生 / 教师 / 管理员 \\
\bottomrule
\end{tabular}
\end{center}

\textbf{服务类：客户端 IUserClientSrv}
\begin{center}
\begin{tabular}{l l l}
\toprule
序号 & 名称 & 方法 \\
\midrule
1 & 登录 & \texttt{User login(User)} \\
2 & 注册 & \texttt{boolean register(User)} \\
3 & 登出 & \texttt{boolean logout(User)} \\
\bottomrule
\end{tabular}
\end{center}

\textbf{服务类：服务器端 IUserServerSrv}
\begin{center}
\begin{tabular}{l l l}
\toprule
序号 & 名称 & 方法 \\
\midrule
1 & 登录校验 & \texttt{User login(User)} \\
2 & 注册 & \texttt{boolean register(User)} \\
3 & 登出 & \texttt{boolean logout(User)} \\
\bottomrule
\end{tabular}
\end{center}

% ========================= 4 其他模块 =========================
\section{其他业务模块设计说明}
学生学籍管理、选课系统、图书馆、商店、银行各模块均采用统一的四层结构与 Message 通信，由对应负责人按 TDD 编写；各模块实体类与数据库表在设计说明书中相应章节补充。

% ========================= 5 公共模块 =========================
\section{公共模块设计说明}
公共工程中同时用于服务器端与客户端的接口、类与资源。用于网络传输的类必须实现 \texttt{java.io.Serializable}，保证两端类一致性。

\textbf{消息类 Message}
\begin{center}
\begin{tabular}{l l l l}
\toprule
序号 & 名称 & 类型 & 备注 \\
\midrule
1 & uid & Long & 消息唯一标识 \\
2 & name & String & 命令名 \\
3 & type & String & 命令 / 数据 \\
4 & statusCode & String & 状态码 \\
5 & data & Object & 传输数据（可序列化） \\
6 & sender & String & 发送者用户名 \\
\bottomrule
\end{tabular}
\end{center}

\textbf{实体类 User}、\textbf{数据库工具类 DbHelper}、\textbf{数据存取类 DAO} 与工具类均置于公共工程中。

% ========================= 6 网络模块 =========================
\section{网络模块设计说明}

\textbf{客户端 IMessageClientSrv}
\begin{center}
\begin{tabular}{l l l}
\toprule
序号 & 名称 & 方法 \\
\midrule
1 & 发送 & \texttt{boolean send(Message)} \\
2 & 接收 & \texttt{boolean receive(Message)} \\
3 & 校验 & \texttt{boolean validate(Message)} \\
\bottomrule
\end{tabular}
\end{center}

\textbf{服务器端 IMessageServerSrv}
\begin{center}
\begin{tabular}{l l l}
\toprule
序号 & 名称 & 方法 \\
\midrule
1 & 发送 & \texttt{boolean send(Message)} \\
2 & 接收 & \texttt{boolean receive(Message)} \\
3 & 校验 & \texttt{boolean validate(Message)} \\
\bottomrule
\end{tabular}
\end{center}

\textbf{Socket 通信}
\begin{itemize}
  \item 客户端连接：\texttt{new Socket(SERVER\_ADDRESS, SERVER\_PORT)}
  \item 服务器端监听：\texttt{new ServerSocket(IConstant.SERVER\_PORT)}
  \item 端口号在公共常量中统一维护，两端保持一致。
\end{itemize}

\textbf{输入输出流（对象流）}
\begin{itemize}
  \item 读取：\texttt{ObjectInputStream.readObject()} → \texttt{Message}
  \item 发送：\texttt{ObjectOutputStream.writeObject(msg)} + \texttt{flush()}
\end{itemize}

% ========================= 7 多线程模块 =========================
\section{多线程模块设计说明}

\textbf{客户端主程序}
\begin{verbatim}
SwingUtilities.invokeLater(new Runnable() {
    public void run() {
        new MainFrame(title).setVisible(true);
    }
});
\end{verbatim}

\textbf{服务器端线程管理}
\begin{center}
\begin{tabular}{l l l}
\toprule
序号 & 名称 & 说明 \\
\midrule
1 & ClientThread & 为每个连接的客户端建立一个线程 \\
2 & ClientThreadMan & 线程管理类，统一管理客户端线程 \\
\bottomrule
\end{tabular}
\end{center}

% ========================= 8 数据库设计 =========================
\section{数据库设计说明}

\subsection{元数据}
数据库 \texttt{vCampus}（MySQL 8.0），主要数据表见下表（其余表由各模块负责人补充）：
\begin{center}
\begin{tabular}{l l l}
\toprule
序号 & 名称 & 说明 \\
\midrule
1 & tblUser & 用户表 \\
2 & tblClass & 班级表 \\
3 & tblXxx & 各业务模块数据表（学籍/选课/图书/商品/账户等） \\
\bottomrule
\end{tabular}
\end{center}

\subsection{表设计：tblUser}
\begin{center}
\begin{tabular}{l l l l l l}
\toprule
字段 & 类型 & 主键 & 默认值 & 约束 & 备注 \\
\midrule
uId & VARCHAR(8) & PRI & NULL & 唯一 & 登录 ID \\
uName & VARCHAR(20) &  & NULL &  & 姓名 \\
uAge & INT &  & NULL & 0--100 & 年龄 \\
uSex & VARCHAR(4) &  & NULL & 男/女 & 性别 \\
uPwd & VARCHAR(32) &  & NULL & 6--16 位 & 密码 \\
uRole & VARCHAR(10) &  & 学生 & 学生/教师/管理员 & 角色 \\
\bottomrule
\end{tabular}
\end{center}

建库与测试数据脚本见 \texttt{sql/vCampus.sql}；各模块表结构由对应负责人按同一命名规范（表名 \texttt{tblXxx}）设计。

% ========================= 附录 =========================
\section*{附录}
本说明书所描述的设计均以仓库中 \texttt{CONTEXT.md}（领域词汇）、\texttt{docs/adr/}（架构决策）、\texttt{docs/project-plan.md}（依赖拓扑与开发规划）为准。

\end{document}
